\clearpage
\section{Security-Betrachtung}
\label{sec:security}

Die Security-Betrachtung konzentriert sich auf die Bluetooth-Kommunikation
zwischen den beiden Raspberry Pis. Die lokale Abstandsauswertung auf Raspberry
Pi 2 bleibt auch dann aktiv, wenn die Funkverbindung gestört oder unterbrochen
wird. Security-Maßnahmen dürfen daher nicht die Fail-Safe-Reaktion verhindern.
Drei Bedrohungen werden durch konkrete Maßnahmen im Design adressiert. Drei
weitere werden bewusst nicht vollständig adressiert, weil dafür zusätzliche
Hardware- oder Plattformmaßnahmen erforderlich wären.

\subsection{Übersicht}

\begin{table}[H]
    \centering
    \small
    \begin{tabularx}{\textwidth}{@{}p{0.8cm}p{3.5cm}YY@{}}
        \toprule
        \textbf{ID} & \textbf{Bedrohung} & \textbf{Mögliche Auswirkung} & \textbf{Maßnahme im Demonstrator} \\
        \midrule
        T1 & Unberechtigtes Gerät sendet Befehle & Falsche Geschwindigkeitsstufe oder unerwünschte Bedienung & Bluetooth-Pairing/Bonding; Pi 2 akzeptiert nur den vorgesehenen Kommunikationspartner; Nachrichteninhalt wird geprüft. \\
        T2 & Alte oder doppelte Nachricht wird erneut verarbeitet & Veralteter Befehl könnte erneut wirksam werden & Sequenznummer je Nachricht; alte und doppelte Nachrichten werden verworfen. \\
        T3 & Bluetooth wird gestört oder fällt aus & Pi 2 erhält keine aktuellen Befehle mehr & Heartbeat und 300-ms-Timeout; bei Kommunikationsverlust Motor AUS. Verfügbarkeit kann dadurch verloren gehen, Safety bleibt erhalten. \\
        \midrule
        T4 & Kompromittierter Raspberry Pi 1 sendet formal gültige, aber manipulierte Befehle & Unerwünschte Geschwindigkeitsstufe, falsche Statusanzeige für den Bediener & \textit{Nicht vollständig adressiert.} Pi 2 entscheidet lokal über die Abschaltung; der Sicherheitsabstand gilt für jede Stufe. \\
        T5 & Manipulation von Software oder Konfiguration auf Pi 2 & Grenzwerte, Timeouts oder Sicherheitsabstände werden verändert oder die Safety-Logik deaktiviert & \textit{Nicht vollständig adressiert.} Konstanten zentral definiert und über Git versioniert; kein Schutz am laufenden Gerät. \\
        T6 & Akustische Störung oder Spoofing der HC-SR04 & Beide Sensoren liefern gleichzeitig falsche Abstände & \textit{Nicht vollständig adressiert.} Plausibilitätsprüfung und zeitversetztes Triggern; gleichartige Manipulation beider Sensoren wird nicht sicher erkannt. \\
        \bottomrule
    \end{tabularx}
    \caption{Security-Bedrohungen und Gegenmaßnahmen}
    \label{tab:security-threats}
\end{table}

\subsection{Adressierte Bedrohungen}

\subsubsection*{T1 -- Unberechtigtes Gerät sendet Befehle}

\textbf{Bedrohung:} Ein fremdes Bluetooth-Gerät baut eine Verbindung zu
Raspberry Pi 2 auf und sendet Steuerbefehle.

\textbf{Angriffsweg:} Ein Angreifer in Funkreichweite gibt sich als
Kommunikationspartner aus oder sendet selbst erzeugte Nachrichten mit
unzulässigem Inhalt.

\textbf{Erkennung und Maßnahme:} Pi 2 akzeptiert nur Verbindungen des
gepaarten bzw.\ gebondeten Raspberry Pi 1. Jede eingehende Nachricht wird vor
der Verwendung auf Nachrichtentyp, Länge und gültigen Wertebereich geprüft.
Nur die Stufen LOW, MEDIUM und HIGH werden als Befehl akzeptiert; unbekannte
Werte führen nicht zu einem Default-Betrieb.

\textbf{Reaktion:}
\begin{center}
    Unbekannter Absender oder ungültiger Nachrichteninhalt $\rightarrow$ Nachricht verwerfen, keine Motorfreigabe
\end{center}

\textbf{Test:} Verbindungsversuch mit einem nicht gepaarten Gerät durchführen
sowie im Testmodus Nachrichten mit falschem Typ, falscher Länge oder
unzulässiger Stufe einspeisen. Die Verbindung bzw.\ Nachricht muss abgelehnt
werden; der Motorzustand darf sich dadurch nicht ändern.

\subsubsection*{T2 -- Wiedereinspielen alter oder doppelter Nachrichten}

\textbf{Bedrohung:} Eine zuvor gültige Nachricht wird erneut verarbeitet und
ein veralteter Befehl wird dadurch wieder wirksam.

\textbf{Angriffsweg:} Aufgezeichnete Nachrichten werden erneut gesendet
(Replay), oder Nachrichten treffen durch Übertragungsfehler doppelt bzw.\ in
falscher Reihenfolge ein.

\textbf{Erkennung und Maßnahme:} Jede Steuerungsnachricht von Pi 1 trägt eine
fortlaufende Sequenznummer. Pi 2 speichert die zuletzt akzeptierte
Sequenznummer und verwirft Nachrichten mit gleicher oder älterer Nummer. Eine
verworfene Nachricht gilt nicht als gültiger Heartbeat und kann den
300-ms-Timeout daher nicht zurücksetzen.

\textbf{Reaktion:}
\begin{center}
    Alte oder doppelte Sequenznummer $\rightarrow$ Nachricht verwerfen; bleibt eine gültige Nachricht aus $\rightarrow$ Motor AUS
\end{center}

\textbf{Test:} Im Testmodus eine bereits verarbeitete Nachricht erneut senden
und Nachrichten mit veralteter Sequenznummer einspeisen. Die Nachrichten
dürfen keinen Befehl auslösen; bei ausschließlich veralteten Nachrichten muss
der Motor nach spätestens 300\,ms abschalten.

\subsubsection*{T3 -- Störung oder Ausfall der Bluetooth-Verbindung}

\textbf{Bedrohung:} Raspberry Pi 2 erhält keine aktuellen Befehle mehr und
könnte mit einem veralteten Befehl weiterarbeiten.

\textbf{Angriffsweg:} Gezielte Funkstörung (Jamming), Unterbrechung der
Verbindung oder Ausfall von Pi 1.

\textbf{Erkennung und Maßnahme:} Pi 1 sendet alle 100\,ms einen Heartbeat mit
dem aktuellen Befehl. Pi 2 überwacht mit einer monotonen Uhr den Zeitpunkt der
letzten gültigen Nachricht. Trifft länger als 300\,ms keine gültige Nachricht
ein, wird ein Kommunikationsfehler gesetzt. Die lokale Abstandsauswertung auf
Pi 2 bleibt dabei unabhängig von der Funkverbindung aktiv.

\textbf{Reaktion:}
\begin{center}
    Kein gültiger Heartbeat innerhalb von 300\,ms $\rightarrow$ Motor AUS
\end{center}

Der Angriff kann damit die Verfügbarkeit des Demonstrators beeinträchtigen,
führt aber nicht zu einem unsicheren Zustand.

\textbf{Test:} Während des Motorbetriebs die Bluetooth-Verbindung trennen oder
das Programm auf Pi 1 anhalten (vgl.\ TC-03). Spätestens nach 300\,ms muss Pi 2
den Motor abschalten.

\subsection{Nicht vollständig adressierte Bedrohungen}

\subsubsection*{T4 -- Kompromittierter Raspberry Pi 1}

\textbf{Bedrohung:} Ein Angreifer erlangt Kontrolle über Raspberry Pi 1 und
sendet Nachrichten, die Typ-, Längen-, Wertebereichs- und Sequenzprüfung
bestehen, inhaltlich aber nicht dem Willen des Bedieners entsprechen.

\textbf{Warum relevant:} Pairing/Bonding und Sequenznummern stellen nur sicher,
dass eine Nachricht vom vorgesehenen Gerät stammt. Sie sagen nichts darüber
aus, ob dieses Gerät vertrauenswürdig arbeitet. Zusätzlich kann ein
manipulierter Pi 1 dem Bediener am LCD einen falschen Systemzustand anzeigen.

\textbf{Warum nicht vollständig adressiert:} Es gibt keine Integritätsprüfung
der Software auf Pi 1 und keine unabhängige Bedienerbestätigung auf Pi 2.

\textbf{Maßnahme im Demonstrator:} Die Abstandsüberwachung läuft lokal auf
Pi 2 und kann durch Befehle von Pi 1 nicht deaktiviert werden. Auch bei
manipulierter Stufe gilt der zugehörige Sicherheitsabstand. Die Auswirkung
beschränkt sich damit im Wesentlichen auf Verfügbarkeit und
Bedienerinformation.

\textbf{Industrielle Lösung:} Secure Boot, gehärtetes Betriebssystem,
Ende-zu-Ende-authentifizierte Nachrichten sowie eine unabhängige Statusanzeige
bzw.\ ein Not-Halt direkt an der Maschine.

\subsubsection*{T5 -- Manipulation von Software oder Konfiguration auf Raspberry Pi 2}

\textbf{Bedrohung:} Durch physischen Zugriff, beispielsweise durch Austausch
der SD-Karte, oder durch Zugriff auf das Betriebssystem werden Programm oder
Parameter von Pi 2 verändert.

\textbf{Warum relevant:} Pi 2 trägt die gesamte Safety-Entscheidung. Eine
Änderung der Sicherheitsabstände, des 300-ms-Timeouts oder der
Plausibilitätsgrenze hebelt die Maßnahmen gegen H1--H3 direkt aus, ohne dass
dies von außen erkennbar ist.

\textbf{Warum nicht vollständig adressiert:} Der Raspberry Pi 4B besitzt im
Projektaufbau keine verifizierte Boot-Kette und keinen Schutz der
Speichermedien. Ein Angreifer mit physischem Zugriff liegt außerhalb des
Projektumfangs (siehe Abschnitt~\ref{subsec:security-abgrenzung}).

\textbf{Maßnahme im Demonstrator:} Safety-relevante Konstanten sind zentral
definiert, Änderungen werden über Git versioniert und vor dem Merge geprüft.
Damit ist der Soll-Stand nachvollziehbar; der Ist-Stand am Gerät wird jedoch
nicht überprüft.

\textbf{Industrielle Lösung:} Secure Boot mit signierter Firmware,
schreibgeschützte bzw.\ verschlüsselte Parameter, Zugriffsschutz am Gehäuse
und Prüfsummen über Safety-Parameter beim Start.

\subsubsection*{T6 -- Akustische Störung oder Spoofing der Ultraschallsensoren}

\textbf{Bedrohung:} Ein Angreifer stört die HC-SR04 gezielt mit Ultraschall
im Bereich von 40\,kHz oder erzeugt künstliche Echos.

\textbf{Warum relevant:} Beide Sensoren arbeiten mit demselben Messprinzip und
derselben Frequenz. Eine externe Quelle kann daher beide gleichzeitig
beeinflussen. Je nach Angriff entstehen zu kleine Messwerte (unnötiges
Abschalten, Verlust der Verfügbarkeit) oder Messwerte, die ein tatsächlich
nahes Objekt verdecken.

\textbf{Warum nicht vollständig adressiert:} Der Quervergleich erkennt nur
Unterschiede zwischen den Sensoren. Werden beide gleichartig manipuliert,
besteht die Plausibilitätsprüfung. Es handelt sich um einen vorsätzlich
herbeigeführten Common-Cause-Fehler (vgl.\ H6).

\textbf{Maßnahme im Demonstrator:} Zeitversetztes Triggern, Gültigkeits- und
Plausibilitätsprüfung sowie Verwendung des kleineren Messwerts. Die
Einschränkung wird ausdrücklich dokumentiert.

\textbf{Industrielle Lösung:} Diversitäre Sensorik, beispielsweise Ultraschall
kombiniert mit Lichtvorhang oder Safety-Laserscanner, codierte Messsignale und
physische Schutzeinrichtungen.

\subsection{Abgrenzung}
\label{subsec:security-abgrenzung}

Der Demonstrator implementiert kein eigenes kryptographisches Protokoll und
keinen vollständigen Schutz gegen einen kompromittierten Raspberry Pi oder
gezieltes Funk-Jamming. Ein Angreifer mit physischem Zugriff auf das Gerät oder
auf das Betriebssystem liegt außerhalb des Projektumfangs. Entscheidend für das
Safety-Konzept ist, dass fehlerhafte oder ausbleibende Kommunikation keine
Motorfreigabe aufrechterhält.
